\section*{Summary of the Base Paper and of the Network Estimators}
\addcontentsline{toc}{section}{Summary of the base paper}
% Budget: 1 page (paper structure v2). Merges the group's summary (A. Lancelotti) with the earlier draft.

\citet{diebold2014network} propose to measure how connected a set of assets is through the variance decomposition of an approximating VAR. The model is estimated and used to forecast each variable $H$ periods ahead; $\theta_{ij}(H)$ is the share of the forecast-error variance of variable $i$ due to shocks in variable $j$. Because shocks are usually correlated, the authors use the generalized variance decomposition of \citet{koop1996impulse} and \citet{pesaran1998generalized}, whose results do not depend on the ordering of the variables; its rows no longer sum to one, so they are normalized. The resulting matrix, augmented with its row and column sums, is the \emph{connectedness table}. Each off-diagonal entry measures how much $j$ affects $i$; the off-diagonal row sum gives how much a variable receives from the others (``from''), the column sum how much it transmits to them (``to''), and their difference whether it is a net transmitter or a net receiver. The average of all off-diagonal entries is the total connectedness of the system. Re-estimating the model over a rolling window shows how connectedness changes over time.

The main conceptual contribution of the paper is that the connectedness table is a weighted, directed network: ``from'' and ``to'' connectedness are the in- and out-degrees of each node, and total connectedness is the mean degree. The two directions also correspond to known systemic-risk measures: what a firm receives resembles exposure measures such as marginal expected shortfall, and what it transmits resembles contribution measures such as $\Delta$CoVaR.

In the application, the authors study the daily realized volatilities of 13 large US financial institutions from May 1999 to April 2010, with a VAR(3), a 12-day horizon and a 100-day rolling window. Total connectedness is 78.3\% in the full sample. Connectedness ``from'' others is tightly distributed (70--82\%), whereas connectedness ``to'' others ranges from 53\% to 106\%, with Citigroup as the largest net transmitter. Over time the index shows two large cycles, around 2000--2003 and during the 2007--2009 crisis, when it reaches its maximum of 89.2\% after the Lehman bankruptcy; the results are robust to the window, the horizon and the ordering of the variables. Our application changes the set of assets (emerging-market country funds), the object (returns instead of volatilities) and the frequency (weekly instead of daily).

Three limitations of the approach are relevant for our work. First, the VAR has $N^2p$ coefficients and is often estimated with few observations. Second, every pair of variables ends up connected, even with a very small weight, so the network never has zero links. Third, the variance decomposition cannot tell whether assets move together because they affect each other or because they are all hit by the same shock. The two estimators below address these limitations.

NETS \citep{barigozzi2019nets} models the panel as a VAR in which both the autoregressive matrices and the precision (or concentration) matrix of the innovations (the inverse of their covariance matrix) are sparse, that is, most of their elements are zero. The dependence structure is then represented by networks. The Granger network is directed: there is a link from $j$ to $i$ when the coefficient on the lags of $j$ in the equation for $i$ is different from zero, so that $j$ has predictive power for $i$. The contemporaneous network is undirected and is built from the partial correlations of the VAR innovations, which measure the relationship between two variables in the same period after controlling for all the others. The two combine into a long-run partial-correlation network. Estimation is carried out jointly by a lasso-type procedure that sets irrelevant links to zero, and, unlike the Diebold--Y{\i}lmaz approach, it requires no identifying assumptions on the shocks. In their application to the volatilities of 90 large US stocks, the authors find that the network captures part of the variability not explained by common volatility factors and improves out-of-sample forecasts. Its main limitation is that it does not model common factors explicitly: when a factor affects the whole panel, the dependence between series is no longer sparse and the central assumption of the model fails.

FNETS \citep{barigozzi2024fnets} addresses this limitation with a factor-adjusted VAR. Each series is decomposed into a common component, driven by a small number of factors in a generalized dynamic factor model \citep{forni2000generalized} and estimated by dynamic principal components, so that factors can affect each series with different lags; and an idiosyncratic component, modeled as a sparse VAR. Since the idiosyncratic component is not observed, its parameters are estimated by an $\ell_1$-regularized Yule--Walker estimator built from the idiosyncratic autocovariances implied by the estimated spectral density. The same three networks as in NETS are then estimated, and they reflect only the dependence that cannot be attributed to the common factors. The method also allows the common and idiosyncratic components to be forecast separately. The authors prove the consistency of the estimators as both the number of series and the sample size grow, also under heavier-than-Gaussian tails, and implement the procedure in the R package \texttt{fnets} \citep{owens2024fnetsR}. For our purpose, FNETS is the key method, because it separates the part of the connectedness among emerging markets that is due to global factors from the part that reflects links between the markets themselves.
